\chapter{Introduction}

% Hook: Electric vehicles need fast charging, long life and safety opration. Understanding battery degradation and side reactoins require monitoring intenral processes during operation.


% Batteries and why impedance matters

The transition to electric mobility and renewable energy storage demands batteries with high energy density, fast charging capability, long operational lifetime and stringent safety requirements. Lithium-ion batteries have dominated the market for portable electronics and are now the technology of choice for electric vehicles, electrical grid stabilization and renewable energy storage. Yet, significant challenges remain in understanding and predicting their degradation mechanisms during operation. Battery aging is a complex interplay of processes occurring at multiple time scales from fast charge transfer reactions at electrode-electrolyte interfaces (microseconds to milliseconds) to slow solid-state diffusion and mechanical degradation (seconds to hours). The understanding and predeiction of a battery, a dynamical system, can be performed from experiment in the time or frequency domain; the latter allowing a much effective way of separating processes at multiple time scale with one simple experiment.

In electrochemistry, the frequency domain identification technique is referred as Electrochemical Impedance Spectroscopy (EIS), excensivly used in battery reserach. By measuring the frequency-dependent response of a battery to small alternating current or voltage perturbations, EIS separates overlapping processes based on their characteristic time constants. A single impedance spectrum encodes information about charge transfer resistance, double-layer capacitance, solid-state diffusion, and interfacial phenomena across a frequency range typically spanning from millihertz to kilohertz. This multi-scale view makes EIS invaluable for battery diagnosis, state-of-health estimation, and mechanistic studies.

% Current state and limitations
However, conventional EIS measurements require the battery to be in a stationary state where the system state (the set of its properties) do not change during the measurement time. For a full spectrum covering six or seven decades of frequency, this requirement translates to measurement times of 30 minutes to several hours. During battery cycling, interrupting the charge or discharge to perform EIS introduces significant overhead time and alters the electrochemical and thermal state of the cell. As a result, operando impedance measurements have been limited to either: narrow frequency bands in the middle to high range (above 1 Hz), which sacrifices the information of transport, or sparse equilibrium measurements at fixed states of charge, which miss the continuous evolution of degradation during long-term cycling. This gap becomes critical when studying phenomena that evolve continuously during cycling: solid-electrolyte interphase growth, lithium plating, particle cracking, or state-dependent kinetics.

% Proposed solution
This thesis addresses this gap by developing a complete methodology for broadband electrochemical impedance spectroscopy in non-stationary conditions. By combining commercial instrumentation with custom automation software and online signal processing, I enable continuous impedance acquisition across seven decades of frequency (10 mHz to 100 kHz) throughout multi-week battery cycling experiments. The methodology removes the stationary-state requirement through multi-frequency perturbation techniques and time-resolved impedance estimation, while solving the practical challenges of data volume (reducing storage requirements by four orders of magnitude through online processing) and multi-channel operation. The open-source software stack makes this capability accessible to other research groups using standard laboratory equipment.

% Explain mulitdisciplinary nature
The development of this methodology required integrating knowledge from multiple disciplines. From electrochemistry, understanding the physical processes at electrode-electrolyte interfaces, charge transfer kinetics, and mass transport phenomena was essential to design appropriate perturbation signals and interpret the resulting impedance spectra. Battery technology knowledge informed the choice of cell geometries, reference electrode designs, and cycling protocols relevant to practical applications. The hardware implementation demanded familiarity with potentiostat operation, waveform generator programming, and high-speed data acquisition systems. Signal processing theory guided the design of multi-sine excitations with optimized crest factors, the implementation of Fourier transform-based impedance estimation, and strategies for managing spectral leakage and noise. System identification concepts—particularly the distinction between stationary and non-stationary systems, transfer function representations, and frequency-domain characterization—provided the theoretical framework for extracting meaningful dynamic models from measured data. Finally, the automation and online processing required advanced programming in Python, including multi-threaded operation, efficient numerical computation with NumPy, and custom hardware interfaces. The parametric analysis of time-varying impedance spectra further demanded expertise in mathematical optimization, including regularized complex nonlinear fitting and model selection criteria. This thesis therefore sits at the intersection of experimental electrochemistry, electrical engineering, digital signal processing, and computational methods—reflecting the inherently multidisciplinary nature of modern battery research.

% Scope statement
The primary focus of this work is the development and validation of the measurement methodology itself. The thesis demonstrates the technique's applicability across multiple electrochemical energy storage systems—commercial batteries, laboratory cells with reference electrodes, Faradaic (intercalation) and non-Faradaic (adsorption) processes—but does not attempt comprehensive mechanistic modeling or full correlation of impedance features to specific degradation modes. Such applications require the foundation that this work provides: a robust, accessible platform for generating the large datasets needed for data-driven battery analysis.


% Thesis outline
This thesis is organized as follows. Part I provides the theoretical framework: Chapter 2 reviews fundamental electrochemistry concepts, Chapter 3 discusses battery technology and characterization methods, Chapter 4 introduces system identification theory for frequency-domain analysis, and Chapter 5 surveys the historical development and state-of-the-art in non-stationary impedance spectroscopy.

Part II describes the methodologies developed in this work: Chapter 7 details the automated data acquisition system, multi-sine signal design, and online impedance calculation, while Chapter 8 presents the analysis approaches for parametric modeling and model selection.

Part III presents results and discussion: Chapter 9 validates the methodology on commercial coin cells throughout their cycling lifetime, Chapters 10-11 demonstrate applications to single-electrode studies (lithium titanium phosphate and electrical double-layer capacitors), and Chapters 12-13 present exploratory measurements on lithium and sodium anodes that reveal methodological challenges for future development. Chapter 14 concludes with a summary of achievements and outlook for future work.